\documentclass[10pt,a4paper]{amsart}
\usepackage[margin=19mm]{geometry}
\usepackage{amsmath,amssymb,mathtools}
\usepackage[scr=rsfs,cal=euler]{mathalfa}
\usepackage{xfrac}
\usepackage[unicode,hidelinks,bookmarks=false]{hyperref}
\newcommand{\ie}{i.e.\ }
\newcommand{\cf}{cf.\ }
\newcommand{\resp}{respectively\ }
\newcommand{\Subsection}{\subsection}
\providecommand{\nobreakdash}{\nobreak\hbox{-}}
\setlength{\emergencystretch}{2em}
\multlinegap0pt
\usepackage{mathtools}
\usepackage{algorithm}\usepackage{algpseudocode}
\usepackage{amssymb}
\usepackage{xcolor}
\newtheorem{lemma}{Lemma}
\newtheorem{proposition}{Proposition}
\newcommand{\C}{\mathbb{C}}
\newcommand{\N}{\mathbb{N}}
\newcommand{\ip}[2]{\langle #1,#2\rangle}
\newcommand{\lin}{\ensuremath{\mathrm{lin}}}
\newcommand{\norm}[1]{\lVert #1\rVert}
\DeclareMathOperator{\spann}{span}
\title{Greedy sampling designs via reduced basis methods:\\optimal recovery in the uniform norm}
\author{Sebastian Neumayer, Kateryna Pozharska, Tino Ullrich}
\date{Source: arXiv:2609.01578v1 (CC BY 4.0)}
\begin{document}
\maketitle

\noindent\emph{Source and licence.} Greedy sampling designs via reduced basis methods:\\optimal recovery in the uniform norm. Version arXiv:2609.01578v1 (CC BY 4.0), distributed by arXiv under the Creative Commons Attribution 4.0 International licence, http://creativecommons.org/licenses/by/4.0/, as recorded in the arXiv licence metadata for this exact version and shown on the abstract page https://arxiv.org/abs/2609.01578v1. Full licence notice: \texttt{licenses/arxiv-2609.01578-CC-BY-4.0.txt}.\par\smallskip

\noindent\textit{Editorial scope.} Counted unit: the article's proposition prop:gm (source0.tex lines 367--372). The article's complete original proof of it is delivered with it (source0.tex lines 374--418, one \texttt{proof} environment of 45 lines). No mathematics is written, repaired or re-derived: the statement and the proof are the article's own lines, in the article's own notation, and no retained line is substituted or re-flowed.  Retained uncounted as the source-local context the proof needs: lines 121--126; lines 351--358.  Nothing was omitted from any retained span, so the package carries no omission record.  No retained line contains a citation, so the wrapper carries no bibliography and no bibliography entry.  No retained line contains a figure environment, an image-inclusion command or drawing code, so no image is delivered and the package declares no asset dependency.  Editorial formatting only, in addition: an AMS \texttt{amsart} preamble that restates the article's own declarations and macros used by the retained lines, namely the environments \texttt{lemma}, \texttt{proposition} on separate counters, together with the standard packages those lines need.  The retained environments print in this document as Lemma 1, Proposition 1 in order of appearance, whereas the article prints the same items with the numbering of its own full document.  The complete native source of the article as distributed in the arXiv e-print for this exact version is bundled unchanged as \texttt{sources/209127/source0.tex}.
\begin{align}
c_n(\mathcal{F})_\infty &\coloneqq \inf_{\substack{\psi\colon\C^n\to B(D)\\ N\colon F \to \C^n \text{ bounded linear}}}
\sup_{f\in \mathcal{F}}\norm{f-\psi(Nf)}_\infty,\\
g_m^{\lin}(\mathcal{F})_\infty &\coloneqq \inf_{\substack{x_1,\dots,x_m\in D\\ \varphi\colon \C^m\to B(D) \text{ linear}}}
\sup_{f\in \mathcal{F}} \bigl\| f-\varphi\bigl(f(x_1),\ldots,f(x_m)\bigr)\bigr\|_\infty \label{eq:glin}
\end{align}

\begin{lemma}[Kernel interpolation is a sampling algorithm]\label{lem:interp}
Let $P=\{x_1,\dots,x_k\}\subset D$ be finite, let $\mathbf K[P]\coloneqq(K(x_i,x_j))_{i,j\le k}$ be the associated Gram matrix with Moore--Penrose inverse $\mathbf K[P]^{+}$, and let $\mathbf f[P]\coloneqq(f(x_j))_{j\le k}$ collect the samples of $f\in\mathcal H_K$.
Then, the $\mathcal H_K$-orthogonal projection onto $V_P \coloneqq \spann\{K(\cdot,x_1),\ldots, K(\cdot,x_k)\}$ is given by
\begin{equation}\label{eq:interp}
    \Pi_{V_P}f=\sum_{i\le k}\bigl(\mathbf K[P]^{+}\mathbf f[P]\bigr)_i K(\cdot,x_i), \qquad f\in\mathcal H_K .
\end{equation}
Thus, $\Pi_{V_P}f$ is linear in the samples, interpolates them, and has minimal $\mathcal H_K$-norm among all functions with these values.
\end{lemma}

\begin{proposition}\label{prop:gm}
For every $m\in\N_0$, it holds that
\begin{equation}\label{eq:chain_prop_gm}
    g_m(B_{\mathcal H_K})_\infty = g_m^{\lin}(B_{\mathcal H_K})_\infty = \tau_m(\mathcal{K})_{\mathcal{H}_K}.
\end{equation}
\end{proposition}

\begin{proof}
Since linear algorithms are a subclass of all algorithms, we have
\begin{equation}\label{eq:gm-triv}
    g_m(B_{\mathcal H_K})_\infty \le g_m^{\lin}(B_{\mathcal H_K})_\infty.
\end{equation}
Below, we prove $g_m^{\lin}(B_{\mathcal H_K})_\infty\le\tau_m(\mathcal{K})_{\mathcal{H}_K}$ and $g_m(B_{\mathcal H_K})_\infty\ge\tau_m(\mathcal{K})_{\mathcal{H}_K}$.
Together with \eqref{eq:gm-triv}, this implies
\begin{equation}
    \tau_m(\mathcal{K})_{\mathcal{H}_K} \le g_m(B_{\mathcal H_K})_\infty \le g_m^{\lin}(B_{\mathcal H_K})_\infty \le \tau_m(\mathcal{K})_{\mathcal{H}_K},
\end{equation}
so that both equalities in \eqref{eq:chain_prop_gm} hold.
 
First, we establish the upper bound for $g_m^{\lin}$.
Fix a set $P=\{x_1,\dots,x_k\}\subset D$ with $k=|P|\le m$.
This is admissible in \eqref{eq:glin}, since repeating a node changes neither $V_P$ nor the resulting algorithm, and by Lemma~\ref{lem:interp} kernel interpolation $f\mapsto\Pi_{V_P}f$ is a linear algorithm using only the samples $\mathbf f[P]$.
Since $f-\Pi_{V_P}f\perp V_P$, we get for $f\in B_{\mathcal H_K}$ and $x\in D$ that
\begin{align}
    |f(x)-\Pi_{V_P}f(x)| &= \bigl|\ip{f-\Pi_{V_P}f}{a_x}_{\mathcal H_K}\bigr|
    = \bigl|\ip{f-\Pi_{V_P}f}{(I-\Pi_{V_P})a_x}_{\mathcal H_K}\bigr| \notag\\
    & \le \norm{f-\Pi_{V_P}f}_{\mathcal H_K} \norm{ (I-\Pi_{V_P})a_x  }_{\mathcal H_K}\le \norm{f}_{\mathcal H_K} \mathrm{Pow}_P(x) \le \mathrm{Pow}_P(x).
\end{align}
Taking the supremum over $x\in D$ and over $f\in B_{\mathcal H_K}$, it holds for every $P$ with $|P|\le m$ that
\begin{equation}
    g_m^{\lin}(B_{\mathcal H_K})_\infty \le
    \sup_{f\in B_{\mathcal H_K}}\norm{f-\Pi_{V_P}f}_\infty \le \norm{\mathrm{Pow}_P}_\infty.
\end{equation}
 
Now, we show the lower bound for $g_m$.
To this end, we use a fooling argument, constructing two admissible functions that produce the same data but are far apart, so that no algorithm can reconstruct both equally well.
Let an arbitrary algorithm based on the samples $f(x_1),\dots,f(x_k)$ at locations $P=\{x_1,\dots,x_k\}$ with $k=|P|\le m$ be given, and let $A\in B(D)$ be its output on the data $(0,\dots,0)$.
We fix $0<\delta<\norm{\mathrm{Pow}_P}_\infty$ and pick $x^*$ with $\mathrm{Pow}_P(x^*)\ge \norm{\mathrm{Pow}_P}_\infty-\delta>0$.
If $\norm{\mathrm{Pow}_P}_\infty=0$ there is nothing to prove.
The function
\begin{equation}
    h \coloneqq \frac{(I-\Pi_{V_P})a_{x^*}}{\mathrm{Pow}_P(x^*)}
\end{equation}
satisfies $\norm{h}_{\mathcal H_K}=1$, $h(x_i)=\ip{h}{a_{x_i}}_{\mathcal H_K}=0$ for all $i\le k$, and $h(x^*)=\ip{h}{a_{x^*}}_{\mathcal H_K}=\mathrm{Pow}_P(x^*)$.
Thus, both $h$ and $-h$ belong to $B_{\mathcal H_K}$ and produce the same data $h(x_i)=-h(x_i)=0$ for $x_i\in P$, so that the
algorithm returns the same $A$ in either case. By the triangle inequality, we have that
\begin{equation}
    \max\bigl\{\norm{h-A}_\infty, \norm{-h-A}_\infty\bigr\} \ge \norm{h}_\infty \ge |h(x^*)| = \mathrm{Pow}_P(x^*)\ge\norm{\mathrm{Pow}_P}_\infty-\delta,
\end{equation}
so that the worst-case error of the recovery algorithm is at least $\norm{\mathrm{Pow}_P}_\infty-\delta$.
Letting $\delta\to0$ and taking the infimum over $P$ gives $g_m(B_{\mathcal H_K})_\infty\ge\tau_m(\mathcal{K})_{\mathcal{H}_K}$.
\end{proof}

\end{document}
